% ------------------------------------------------------------
% Packages
% ------------------------------------------------------------
\usepackage{amsmath,amsthm,amssymb,amsbsy}
\usepackage{xcolor}
\usepackage{graphicx}
\graphicspath{{./figs/}}
\usepackage{subcaption}
\usepackage{wrapfig}
\usepackage{booktabs}
\usepackage{multirow}
\usepackage{array}
\usepackage{enumitem}
\usepackage{comment}
\usepackage{cleveref} % load after hyperref
\crefname{equation}{Eq.}{Eqs.}
\Crefname{equation}{Eq.}{Eqs.}
\usepackage{algorithm}
\usepackage{algpseudocode}
\usepackage{listings}
\usepackage{pgfplots}
\pgfplotsset{compat=1.18}
\usetikzlibrary{positioning,calc}
\newcommand{\appref}[1]{\hyperref[#1]{Appendix~\ref*{#1}}}

% ------------------------------------------------------------
% Code listings: \begin{lstlisting}[language=Python] ... \end{lstlisting}
% ------------------------------------------------------------
\definecolor{codegray}{RGB}{110,118,129}
\definecolor{codekw}{RGB}{0,79,163}
\definecolor{codestr}{RGB}{158,64,0}
\lstset{
  basicstyle=\ttfamily\footnotesize,
  keywordstyle=\color{codekw}\bfseries,
  stringstyle=\color{codestr},
  commentstyle=\color{codegray}\itshape,
  numbers=left, numberstyle=\tiny\color{codegray}, numbersep=8pt,
  backgroundcolor=\color{gray!4},
  frame=single, rulecolor=\color{gray!35}, framesep=6pt, xleftmargin=16pt, xrightmargin=6.4pt,
  breaklines=true, columns=fullflexible, keepspaces=true, showstringspaces=false,
  aboveskip=10pt, belowskip=6pt,
}

% ------------------------------------------------------------
% Plots: \begin{axis}[arxiv plot] ... ; cycle list matches the template colors
% ------------------------------------------------------------
\pgfplotscreateplotcyclelist{arxiv}{
  {ArxivAccent, very thick, mark=none},
  {ArxivLinkBlue, very thick, mark=none},
  {takeawaygreen, very thick, mark=none, dashed},
  {gray, thick, mark=none, densely dotted},
}
\pgfplotsset{
  arxiv plot/.style={
    width=\linewidth, height=0.56\linewidth,
    cycle list name=arxiv,
    axis lines=left, axis line style={gray!60},
    tick style={gray!60}, tick label style={font=\scriptsize\sffamily},
    label style={font=\footnotesize\sffamily},
    legend style={font=\scriptsize\sffamily, draw=none, fill=none},
    legend cell align=left,
    grid=major, major grid style={gray!15},
    every axis plot/.append style={line join=round},
  },
}

% ------------------------------------------------------------
% Theorems
% ------------------------------------------------------------
\theoremstyle{plain}
\newtheorem{theorem}{Theorem}
\newtheorem{lemma}{Lemma}
\newtheorem{prop}{Proposition}
\newtheorem{cor}{Corollary}
\theoremstyle{definition}
\newtheorem{definition}{Definition}
\theoremstyle{remark}
\newtheorem{remark}{Remark}

% ------------------------------------------------------------
% Callout boxes
% Usage: \researchquestions{\textbf{Q1.}\enspace First?\par\smallskip \textbf{Q2.}\enspace Second?}
%        \takeaway{Main finding.}
%        \begin{promptbox} ... \end{promptbox}
%        \begin{judgepromptbox}{Title} ... \end{judgepromptbox}
% Optional titles: \researchquestions[Research question]{...}, \takeaway[Takeaways]{...}
% ------------------------------------------------------------
\usepackage[most]{tcolorbox}
\usepackage{fontawesome5}
\definecolor{researchblue}{RGB}{62,100,125}
\definecolor{takeawaygreen}{RGB}{62,112,84}
\definecolor{promptpurple}{RGB}{112,82,140}

% \begin{researchbox}[tcolorbox options]{color}
\newtcolorbox{researchbox}[2][]{
  enhanced,
  breakable,
  colback=#2!5,
  frame hidden,
  overlay unbroken and first={
    \begin{tcbclipinterior}
      \fill[#2] (frame.south west) rectangle ([xshift=2.5pt]frame.north west);
    \end{tcbclipinterior}
  },
  overlay middle and last={
    \begin{tcbclipinterior}
      \fill[#2] (frame.south west) rectangle ([xshift=2.5pt]frame.north west);
    \end{tcbclipinterior}
  },
  boxrule=0pt,
  arc=4pt,
  left=9pt, right=9pt, top=7pt, bottom=7pt,
  boxsep=0pt,
  before skip=10pt, after skip=10pt,
  #1,
}
\newtcolorbox{promptbox}{
  enhanced, breakable,
  colback=gray!4, colframe=gray!35, boxrule=0.4pt, arc=3pt,
  left=8pt, right=8pt, top=6pt, bottom=6pt,
  fontupper=\small,
  before skip=10pt, after skip=10pt,
}
\newenvironment{judgepromptbox}[1]{%
  \begin{researchbox}{promptpurple}
    {\color{promptpurple}\faGavel\enspace\sffamily\bfseries #1}\par\smallskip
    \small
}{\end{researchbox}}
\newcommand{\researchquestions}[2][Research questions]{%
  \begin{researchbox}[unbreakable]{researchblue}
    {\color{researchblue}\faQuestionCircle\enspace\sffamily\bfseries #1}\par\smallskip
    #2
  \end{researchbox}%
}
\newcommand{\takeaway}[2][Takeaway]{%
  \begin{researchbox}[unbreakable]{takeawaygreen}
    {\color{takeawaygreen}\faLightbulb\enspace\sffamily\bfseries #1}\par\smallskip
    #2
  \end{researchbox}%
}

% fields
\renewcommand{\O}{\mathbb{O}}
\newcommand{\R}{\mathbb{R}}
\newcommand{\V}{\mathbb{V}}
\def \real    { \mathbb{R} }
\def \reals   { \mathbb{R} }
\def \Real {\mathbb{R}}
\newcommand{\C}{\mathbb{C}}
\def \complex{\mathbb{C}}
\def \Complex {\mathbb{C}}
\newcommand{\Z}{\mathbb{Z}}
\newcommand{\N}{\mathbb{N}}

\newcommand{\la}{\left\langle}
\newcommand{\ra}{\right\rangle}
\newcommand{\e}{
\begin{equation}}
  \newcommand{\ee}{
\end{equation}}
\newcommand{\en}{
\begin{equation*}}
  \newcommand{\een}{
\end{equation*}}
\newcommand{\eqn}{
\begin{eqnarray}}
  \newcommand{\eeqn}{
\end{eqnarray}}
\newcommand{\bmat}{
\begin{bmatrix}}
  \newcommand{\emat}{
\end{bmatrix}}
% complex numbers
\renewcommand{\Re}[1]{\operatorname{Re}\left\{#1\right\}}
\renewcommand{\Im}[1]{\operatorname{Im}\left\{#1\right\}}
\newcommand{\conj}[1]{\mkern 1.5mu\overline{\mkern-1.5mu#1\mkern-1.5mu}\mkern 1.5mu}

\DeclareMathAlphabet\mathbfcal{OMS}{cmsy}{b}{n}
% probability and stat
\renewcommand{\P}[1]{\operatorname{\mathbb{P}}\left(#1\right)}
\newcommand{\E}{\operatorname{\mathbb{E}}}
\newcommand{\dif}{\operatorname{d}}
\newcommand{\var}{\operatorname{var}}

% calculus
\renewcommand{\d}[1]{\,\mathrm{d}#1}

% constants (written in roman, if wanted)
%\newcommand{\e}{\mathrm{e}}
\renewcommand{\j}{\mathrm{j}}

\newcommand{\mb}{\mathbf}
\newcommand{\mc}{\mathcal}
\newcommand{\mf}{\mathfrak}
\newcommand{\bb}{\mathbb}
\newcommand{\mrm}{\mathrm}
\newcommand{\mtt}{\mathtt}
\def \vecword  {\operatorname*{vec}}
% linear algebra
%   vector notation
\newcommand{\vct}[1]{\boldsymbol{#1}}
%   matrices
\newcommand{\mtx}[1]{\boldsymbol{#1}}
%   block vector
\newcommand{\bvct}[1]{\mathbf{#1}}
%   block matrix
\newcommand{\bmtx}[1]{\mathbf{#1}}
%  inner products
\newcommand{\inner}[1]{\left<#1\right>}
\newcommand{\<}{\langle}
\renewcommand{\>}{\rangle}
%   transpose, Hermitian, pseudo-inverse
\renewcommand{\H}{\mathrm{H}}
\newcommand{\T}{\mathrm{T}}
\newcommand{\pinv}{\dagger}
%  fundamental subspaces
\newcommand{\Null}{\operatorname{Null}}
\newcommand{\Range}{\operatorname{Range}}
\newcommand{\Span}{\operatorname{Span}}
%  operators
\newcommand{\trace}{\operatorname{trace}}
\newcommand{\conv}{\operatorname{conv}}
\newcommand{\rank}{\operatorname{rank}}
\newcommand{\diag}{\operatorname{diag}}
\newcommand{\toep}{\operatorname{Toep}}
\newcommand{\SDP}{\operatorname{SDP}}
\newcommand{\PSD}{\operatorname{PSD}}
\newcommand{\far}{\operatorname{far}}
\newcommand{\near}{\operatorname{near}}
\newcommand{\sign}{\operatorname{sign}}
\newcommand{\sgn}{\operatorname{sgn}}
\newcommand{\Sign}{\operatorname{Sgn}}
\newcommand{\dist}{\operatorname{dist}}
\def \lg        {\langle}
\def \rg        {\rangle}
\def \vec       {\operatorname*{vec}}
%
\newcommand{\dimension}{\operatorname{dim}}

% sets and topology
\newcommand{\set}[1]{\mathbb{#1}}
\newcommand{\closure}{\operatorname{cl}}  % closure
\newcommand{\interior}{\operatorname{int}}
\newcommand{\boundary}{\operatorname{bd}}
\newcommand{\diameter}{\operatorname{diam}}

% functional analysis
\newcommand{\domain}{\operatorname{dom}}
\newcommand{\epigraph}{\operatorname{epi}}
\newcommand{\hypograph}{\operatorname{hypo}}
\newcommand{\linop}[1]{\mathscr{#1}}  % general linear operator

% optimization
\DeclareMathOperator*{\minimize}{\text{minimize}}
\DeclareMathOperator*{\maximize}{\text{maximize}}
\DeclareMathOperator*{\argmin}{\text{arg~min}}
\DeclareMathOperator*{\argmax}{\text{arg~max}}
\def \st {\operatorname*{s.t.\ }}
% other
%\newcommand{\sgn}{\text{sgn}}
\newcommand{\sinc}{\text{sinc}}
\newcommand{\floor}[1]{\left\lfloor #1 \right\rfloor}
\newcommand{\ceil}[1]{\left\lceil #1 \right\rceil}
\newcommand{\eps}{\epsilon}

\newcommand{\wh}{\widehat}
\newcommand{\wc}{\widecheck}
\newcommand{\wt}{\widetilde}
\newcommand{\ol}{\overline}

\newcommand{\norm}[2]{\left\| #1 \right\|_{#2}}
\newcommand{\nrm}[1]{\left\Vert#1\right\Vert}
\newcommand{\abs}[1]{\left| #1 \right|}
\newcommand{\bracket}[1]{\left( #1 \right)}
\newcommand{\parans}[1]{\left(#1\right)}
\newcommand{\innerprod}[2]{\left\langle #1,  #2 \right\rangle}
\newcommand{\prob}[1]{\bb P\left[ #1 \right]}
\newcommand{\expect}[1]{\bb E\left[ #1 \right]}
\newcommand{\function}[2]{#1 \left(#2\right)}
\newcommand{\integral}[4]{\int_{#1}^{#2}\; #3\; #4}
\newcommand{\iter}[1]{^{(#1)}}

%--------------------------------------------------------------------------
\newcommand{\calA}{\mathcal{A}}
\newcommand{\calB}{\mathcal{B}}
\newcommand{\calC}{\mathcal{C}}
\newcommand{\calD}{\mathcal{D}}
\newcommand{\calE}{\mathcal{E}}
\newcommand{\calF}{\mathcal{F}}
\newcommand{\calG}{\mathcal{G}}
\newcommand{\calH}{\mathcal{H}}
\newcommand{\calI}{\mathcal{I}}
\newcommand{\calJ}{\mathcal{J}}
\newcommand{\calK}{\mathcal{K}}
\newcommand{\calL}{\mathcal{L}}
\newcommand{\calM}{\mathcal{M}}
\newcommand{\calN}{\mathcal{N}}
\newcommand{\calO}{\mathcal{O}}
\newcommand{\calP}{\mathcal{P}}
\newcommand{\calQ}{\mathcal{Q}}
\newcommand{\calR}{\mathcal{R}}
\newcommand{\calS}{\mathcal{S}}
\newcommand{\calT}{\mathcal{T}}
\newcommand{\calU}{\mathcal{U}}
\newcommand{\calV}{\mathcal{V}}
\newcommand{\calW}{\mathcal{W}}
\newcommand{\calX}{\mathcal{X}}
\newcommand{\calY}{\mathcal{Y}}
\newcommand{\calZ}{\mathcal{Z}}

\newcommand{\va}{\vct{a}}
\newcommand{\vb}{\vct{b}}
\newcommand{\vc}{\vct{c}}
\newcommand{\vd}{\vct{d}}
\newcommand{\ve}{\vct{e}}
\newcommand{\vf}{\vct{f}}
\newcommand{\vg}{\vct{g}}
\newcommand{\vh}{\vct{h}}
\newcommand{\vi}{\vct{i}}
\newcommand{\vj}{\vct{j}}
\newcommand{\vk}{\vct{k}}
\newcommand{\vl}{\vct{l}}
\newcommand{\vm}{\vct{m}}
\newcommand{\vn}{\vct{n}}
\newcommand{\vo}{\vct{o}}
\newcommand{\vp}{\vct{p}}
\newcommand{\vq}{\vct{q}}
\newcommand{\vr}{\vct{r}}
\newcommand{\vs}{\vct{s}}
\newcommand{\vt}{\vct{t}}
\newcommand{\vu}{\vct{u}}
\newcommand{\vv}{\vct{v}}
\newcommand{\vw}{\vct{w}}
\newcommand{\vx}{\vct{x}}
\newcommand{\vy}{\vct{y}}
\newcommand{\vz}{\vct{z}}
%
\newcommand{\valpha}{\vct{\alpha}}
\newcommand{\vbeta}{\vct{\beta}}
\newcommand{\vtheta}{\vct{\theta}}
\newcommand{\vepsilon}{\vct{\epsilon}}
\newcommand{\veps}{\vct{\epsilon}}
\newcommand{\vvarepsilon}{\vct{\varepsilon}}
\newcommand{\vgamma}{\vct{\gamma}}
\newcommand{\vlambda}{\vct{\lambda}}
\newcommand{\vnu}{\vct{\nu}}
\newcommand{\vphi}{\vct{\phi}}
\newcommand{\vpsi}{\vct{\psi}}
\newcommand{\vdelta}{\vct{\delta}}
\newcommand{\veta}{\vct{\eta}}
\newcommand{\vxi}{\vct{\xi}}
\newcommand{\vmu}{\vct{\mu}}
\newcommand{\vrho}{\vct{\rho}}
%
\newcommand{\vzero}{\vct{0}}
\newcommand{\vone}{\vct{1}}

\newcommand{\mA}{\mtx{A}}
\newcommand{\mB}{\mtx{B}}
\newcommand{\mC}{\mtx{C}}
\newcommand{\mD}{\mtx{D}}
\newcommand{\mE}{\mtx{E}}
\newcommand{\mF}{\mtx{F}}
\newcommand{\mG}{\mtx{G}}
\newcommand{\mH}{\mtx{H}}
\newcommand{\mJ}{\mtx{J}}
\newcommand{\mK}{\mtx{K}}
\newcommand{\mL}{\mtx{L}}
\newcommand{\mM}{\mtx{M}}
\newcommand{\mN}{\mtx{N}}
\newcommand{\mO}{\mtx{O}}
\newcommand{\mP}{\mtx{P}}
\newcommand{\mQ}{\mtx{Q}}
\newcommand{\mR}{\mtx{R}}
\newcommand{\mS}{\mtx{S}}
\newcommand{\mT}{\mtx{T}}
\newcommand{\mU}{\mtx{U}}
\newcommand{\mV}{\mtx{V}}
\newcommand{\mW}{\mtx{W}}
\newcommand{\mX}{\mtx{X}}
\newcommand{\mY}{\mtx{Y}}
\newcommand{\mZ}{\mtx{Z}}
%
\newcommand{\mGamma}{\mtx{\Gamma}}
\newcommand{\mDelta}{\mtx{\Delta}}
\newcommand{\mLambda}{\mtx{\Lambda}}
\newcommand{\mOmega}{\mtx{\Omega}}
\newcommand{\mPhi}{\mtx{\Phi}}
\newcommand{\mPsi}{\mtx{\Psi}}
\newcommand{\mSigma}{\mtx{\Sigma}}
\newcommand{\mTheta}{\mtx{\Theta}}
\newcommand{\mUpsilon}{\mtx{\Upsilon}}
%
\newcommand{\mId}{{\bf I}}
\newcommand{\mEx}{{\bf J}}
\newcommand{\mzero}{{\bf 0}}
\newcommand{\mone}{{\bf 1}}
%
\newcommand{\bva}{\bvct{a}}
\newcommand{\bvh}{\bvct{h}}
\newcommand{\bvw}{\bvct{w}}
\newcommand{\bvx}{\bvct{x}}
\newcommand{\bvy}{\bvct{y}}

\newcommand{\bmM}{\bmtx{M}}
\newcommand{\bmH}{\bmtx{H}}
\newcommand{\bmX}{\bmtx{X}}

\newcommand{\setA}{\set{A}}
\newcommand{\setB}{\set{B}}
\newcommand{\setC}{\set{C}}
\newcommand{\setD}{\set{D}}
\newcommand{\setE}{\set{E}}
\newcommand{\setF}{\set{F}}
\newcommand{\setG}{\set{G}}
\newcommand{\setH}{\set{H}}
\newcommand{\setI}{\set{I}}
\newcommand{\setJ}{\set{J}}
\newcommand{\setK}{\set{K}}
\newcommand{\setL}{\set{L}}
\newcommand{\setM}{\set{M}}
\newcommand{\setN}{\set{N}}
\newcommand{\setO}{\set{O}}
\newcommand{\setP}{\set{P}}
\newcommand{\setQ}{\set{Q}}
\newcommand{\setR}{\set{R}}
\newcommand{\setS}{\set{S}}
\newcommand{\setT}{\set{T}}
\newcommand{\setU}{\set{U}}
\newcommand{\setV}{\set{V}}
\newcommand{\setW}{\set{W}}
\newcommand{\setX}{\set{X}}
\newcommand{\setY}{\set{Y}}
\newcommand{\setZ}{\set{Z}}

\renewcommand{\ss}{{\vspace*{-1mm}}}

\setcounter{MaxMatrixCols}{20}

\newlength{\imgwidth}
\setlength{\imgwidth}{3.125in}

% Review comments: comment out \showcommentstrue to hide them and unhighlight revisions.
\newif\ifshowcomments
\showcommentstrue
\ifshowcomments
  \newcommand{\zz}[1]{\textcolor{red}{\bf [{\em ZZ:} #1]}}
  \newcommand{\xz}[1]{\textcolor{purple}{\bf [{\em XZ:} #1]}}
  \newcommand{\revise}[1]{\textcolor{blue}{#1}}
  \newcommand{\edit}[1]{\textcolor{red}{#1}}
\else
  \newcommand{\zz}[1]{}
  \newcommand{\xz}[1]{}
  \newcommand{\revise}[1]{#1}
  \newcommand{\edit}[1]{#1}
\fi
\newcommand{\zzrevise}[1]{\revise{#1}}

\setlist[enumerate]{leftmargin=*}
\setlist[itemize]{leftmargin=*}
